\documentclass[10pt]{beamer}

\usetheme{metropolis}
\setbeamertemplate{navigation symbols}{}
\setbeamercolor{background canvas}{bg=white}
\metroset{progressbar=none}

\title{GNN Neural Surrogates}
\subtitle{Why and how the current papers use program graphs}
\author{}
\institute{}
\date{}

\begin{document}

\maketitle

\begin{frame}{Research Context}
  \begin{block}{Working Idea}
    Convert a program into one or more graphs, pass those graphs through a GNN, and reason over the learned vectors instead of reasoning directly over the original program.
  \end{block}

  \begin{itemize}
    \item The learned vectors act as a \textbf{neural surrogate} for program structure and behavior.
    \item Possible downstream uses: cloud-side analysis without exposing raw code, code transformation, bug localization, vulnerability detection, similarity search, and tool selection.
    \item For now, the useful question is not ``which paper gets the best score?'' but ``why do they need the surrogate, and how do they build it?''
  \end{itemize}
\end{frame}

\begin{frame}{Overview}
  \tableofcontents
\end{frame}

\section{Reading Lens}

\begin{frame}{Common Pipeline}
  \begin{enumerate}
    \item Choose a program view: source code, IR, CFG, DFG, CPG, execution trace, or query/target graph.
    \item Convert that view into a graph: nodes represent tokens, statements, variables, instructions, methods, or basic blocks.
    \item Add typed edges: syntax, control flow, data flow, dependence, call, token order, or trace relation.
    \item Run a GNN: message passing turns local graph neighborhoods into vector representations.
    \item Read out vectors: classify a graph, score nodes, compare graph pairs, retrieve matches, or rank tools.
  \end{enumerate}
\end{frame}

\begin{frame}{How We Should Read Each Paper}
  \begin{itemize}
    \item \textbf{Why:} what original program-analysis task is too hard, too slow, too manual, or too sensitive to do directly?
    \item \textbf{Graph:} what program information is preserved in the surrogate?
    \item \textbf{GNN role:} what kind of message passing or embedding is used?
    \item \textbf{Vector use:} what is done after the vectors are produced?
    \item \textbf{Research lesson:} what design choice is relevant for our neural-surrogate direction?
  \end{itemize}
\end{frame}

\begin{frame}{Paper Categories}
  \footnotesize
  \setlength{\tabcolsep}{2pt}
  \begin{tabular}{p{0.24\linewidth}p{0.25\linewidth}p{0.38\linewidth}}
    \textbf{Category} & \textbf{Papers} & \textbf{Surrogate Question} \\
    \hline
    Foundation & Allamanis & Why graph representations make sense for programs \\
    Vulnerability detection & Devign, GraphEye & Can vectors summarize vulnerable program semantics? \\
    Trace-based security analysis & Silent buffer overflows & Can vectors replace source-code access when only execution traces exist? \\
    Bug-pattern search & SICode, NeuroMatch & Can vector arithmetic approximate subgraph reasoning? \\
    Program similarity & funcGNN & Can vectors approximate expensive graph similarity? \\
    Verification support & Graves & Can vectors predict which formal tool should analyze a program? \\
  \end{tabular}
\end{frame}

\section{Foundation}

\begin{frame}{Allamanis: GNNs In Program Analysis}
  \begin{columns}[T,onlytextwidth]
    \column{0.48\linewidth}
    \textbf{Why}
    \begin{itemize}
      \item Programs have structure that token sequences flatten away.
      \item Static analyses expose useful relations, but exact analyses can be brittle, incomplete, or task-specific.
      \item GNNs give a way to learn from program structure while keeping relational information explicit.
    \end{itemize}

    \column{0.48\linewidth}
    \textbf{How}
    \begin{itemize}
      \item Represent program elements as graph nodes.
      \item Connect nodes using syntactic, control-flow, data-flow, and variable-use relations.
      \item Use message passing to aggregate local context into node or graph vectors.
    \end{itemize}
  \end{columns}

  \vspace{0.25cm}
  \textbf{Use for us:} vocabulary and motivation for treating the graph embedding as a program surrogate.
\end{frame}

\section{Vulnerability Detection}

\begin{frame}{Devign}
  \small
  \begin{columns}[T,onlytextwidth]
    \column{0.48\linewidth}
    \textbf{Why}
    \begin{itemize}
      \item Vulnerability detection needs semantic clues that manual patterns or token sequences often miss.
      \item The goal is to classify vulnerable functions in real C projects.
      \item The surrogate should preserve enough syntax, control, and data dependency information to support a security decision.
    \end{itemize}

    \column{0.48\linewidth}
    \textbf{How}
    \begin{itemize}
      \item Build a composite code graph for each function.
      \item Include AST-like structure, control-flow edges, data-flow edges, and token-order edges.
      \item Run gated graph recurrent message passing, then use a convolution/readout module for graph-level classification.
    \end{itemize}
  \end{columns}

  \vspace{0.25cm}
  \textbf{Use for us:} ``function $\rightarrow$ semantic graph $\rightarrow$ vector $\rightarrow$ vulnerability label.''
\end{frame}

\begin{frame}{GraphEye}
  \begin{columns}[T,onlytextwidth]
    \column{0.48\linewidth}
    \textbf{Why}
    \begin{itemize}
      \item Manual code review does not scale for C/C++ vulnerability detection.
      \item The authors assume vulnerable and non-vulnerable functions differ in their code property graph structure.
      \item They want a learned graph classifier that focuses on the important parts of the program graph.
    \end{itemize}

    \column{0.48\linewidth}
    \textbf{How}
    \begin{itemize}
      \item Convert source code into a code property graph.
      \item Vectorize that graph with VecCPG to encode syntax and semantic features.
      \item Use a graph attention model, GcGAT, to classify the function graph.
    \end{itemize}
  \end{columns}

  \vspace{0.25cm}
  \textbf{Use for us:} attention is useful when the surrogate graph is large and only some neighborhoods matter.
\end{frame}

\section{Trace-Based Analysis}

\begin{frame}{Silent Buffer Overflows}
  \begin{columns}[T,onlytextwidth]
    \column{0.48\linewidth}
    \textbf{Why}
    \begin{itemize}
      \item Some buffer overflows produce no obvious crash or visible symptom.
      \item Source code may be unavailable, so source-level graph construction is not always possible.
      \item The task becomes reasoning from execution traces instead of raw source.
    \end{itemize}

    \column{0.48\linewidth}
    \textbf{How}
    \begin{itemize}
      \item Extract a DFG+ graph from an execution trace.
      \item Use typed data-flow and trace relations to represent how values move during execution.
      \item Train a modified relational GCN to detect overflow-relevant behavior from the graph.
    \end{itemize}
  \end{columns}

  \vspace{0.25cm}
  \textbf{Use for us:} the surrogate does not have to come from source code; dynamic traces can also become graphs and vectors.
\end{frame}

\section{Bug Pattern And Subgraph Reasoning}

\begin{frame}{SICode}
  \begin{columns}[T,onlytextwidth]
    \column{0.48\linewidth}
    \textbf{Why}
    \begin{itemize}
      \item Given a known buggy snippet, we want to find similar hidden buggy structures in a large target project.
      \item Exact subgraph isomorphism is expensive.
      \item Whole-function similarity can miss cases where the bug is only a small substructure inside a different-looking function.
    \end{itemize}

    \column{0.48\linewidth}
    \textbf{How}
    \begin{itemize}
      \item Slice or extract a query graph around the known buggy code.
      \item Convert target code into featured code graphs.
      \item Train a graph embedding model so subgraph-isomorphism relations can be judged by comparing embedding vectors.
    \end{itemize}
  \end{columns}

  \vspace{0.25cm}
  \textbf{Use for us:} a direct example of replacing an expensive symbolic graph problem with vector-space reasoning.
\end{frame}

\begin{frame}{NeuroMatch}
  \begin{columns}[T,onlytextwidth]
    \column{0.48\linewidth}
    \textbf{Why}
    \begin{itemize}
      \item Subgraph matching is important but NP-complete.
      \item Large query and target graphs make exact combinatorial matching impractical.
      \item The authors want approximate matching directly in embedding space.
    \end{itemize}

    \column{0.48\linewidth}
    \textbf{How}
    \begin{itemize}
      \item Decompose query and target graphs into smaller neighborhoods.
      \item Embed those subgraphs with a GNN.
      \item Train the embedding space to respect subgraph-relation geometry, then retrieve matches by vector comparison.
    \end{itemize}
  \end{columns}

  \vspace{0.25cm}
  \textbf{Use for us:} method background for SICode-style bug retrieval and any future vectorized pattern-matching task.
\end{frame}

\section{Similarity And Tool Selection}

\begin{frame}{funcGNN}
  \begin{columns}[T,onlytextwidth]
    \column{0.48\linewidth}
    \textbf{Why}
    \begin{itemize}
      \item Program similarity supports clone detection, plagiarism detection, code search, and refactoring.
      \item CFG-based graph edit distance captures structure but is computationally expensive.
      \item The goal is to learn a cheaper vector surrogate for similarity.
    \end{itemize}

    \column{0.48\linewidth}
    \textbf{How}
    \begin{itemize}
      \item Convert programs into control-flow graphs.
      \item Use a GNN/Siamese-style setup to embed two program graphs.
      \item Compare graph embeddings to estimate similarity or approximate graph edit distance.
    \end{itemize}
  \end{columns}

  \vspace{0.25cm}
  \textbf{Use for us:} a close analogy for code transformation or migration, where vector-space similarity could guide alignment between languages.
\end{frame}

\begin{frame}{Graves}
  \begin{columns}[T,onlytextwidth]
    \column{0.48\linewidth}
    \textbf{Why}
    \begin{itemize}
      \item Software verification has many tools and algorithms, and no single tool is best on every program.
      \item Choosing the verifier manually requires expertise.
      \item The goal is not to verify directly, but to predict which verifier will work well.
    \end{itemize}

    \column{0.48\linewidth}
    \textbf{How}
    \begin{itemize}
      \item Build a graph representation of the program.
      \item Feed that graph to a GNN or graph-attention model.
      \item Predict a score for each verifier and use the scores to choose or order tools.
    \end{itemize}
  \end{columns}

  \vspace{0.25cm}
  \textbf{Use for us:} a surrogate can be used for orchestration: it can decide what analysis to run, not only output the final answer.
\end{frame}

\section{Synthesis}

\begin{frame}{What The Papers Teach Us}
  \begin{itemize}
    \item Program-to-vector surrogates are already useful for classification, retrieval, similarity, and algorithm selection.
    \item The graph construction is as important as the neural model because it decides what information survives.
    \item Source graphs are good for static code reasoning; trace graphs are good when runtime behavior or source privacy matters.
    \item Graph-pair embeddings are especially relevant for transformation, migration, and pattern matching.
    \item Node-level embeddings are the right direction for localization tasks such as ranking likely buggy statements or variables.
  \end{itemize}
\end{frame}

\begin{frame}{Relation To Our Research Ideas}
  \begin{itemize}
    \item \textbf{Cloud-side private analysis:} closest prior idea is replacing source code with vectors, but these papers do not prove privacy; they mostly prove task usefulness.
    \item \textbf{Code transformation:} funcGNN and subgraph matching suggest ways to align structures, but the current set does not directly solve Java-to-COBOL translation.
    \item \textbf{Error detection and localization:} Devign and GraphEye support graph-level detection; a node-scoring version would need node-level supervision or an attribution mechanism.
    \item \textbf{Tool assistance:} Graves suggests using GNN vectors to choose downstream analyzers instead of replacing them.
  \end{itemize}
\end{frame}

\begin{frame}{Open Design Questions For Us}
  \begin{itemize}
    \item What graph view should be the surrogate: AST, CFG, DFG, CPG, IR graph, trace graph, or a bundle of graphs?
    \item Should the surrogate be graph-level, node-level, edge-level, or graph-pair based?
    \item How much information must the vector preserve for the task, and how much should it intentionally hide for privacy?
    \item Can the vector representation support multiple tasks, or should each task get a specialized surrogate?
    \item For transformation, do we need embeddings alone, or embeddings plus a decoder/generator?
  \end{itemize}
\end{frame}

\begin{frame}{Suggested Reading Order}
  \begin{enumerate}
    \item Start with Allamanis for the general program-graph motivation.
    \item Read Devign and GraphEye for function-level vulnerability surrogates.
    \item Read the silent-buffer-overflow paper to see trace-based surrogates without source-code access.
    \item Read SICode and NeuroMatch for vectorized subgraph reasoning.
    \item Read funcGNN for graph-similarity surrogates.
    \item Read Graves for surrogate-driven tool selection.
  \end{enumerate}
\end{frame}

\begin{frame}{References In This Folder}
  \scriptsize
  \begin{itemize}
    \item Allamanis, ``Graph Neural Networks in Program Analysis'' / \texttt{allamanis2021graph.pdf}
    \item Zhou et al., ``Devign'' / \texttt{1909.03496v1.pdf}
    \item Zhou et al., ``GraphEye'' / \texttt{2202.02501v1.pdf}
    \item Wang et al., ``Spotting Silent Buffer Overflows...'' / \texttt{2102.10452v1.pdf}
    \item Gong et al., ``SICode'' / \texttt{3643916.3646556.pdf}
    \item Ying et al., ``Neural Subgraph Matching'' / \texttt{2007.03092v2.pdf}
    \item Nair et al., ``funcGNN'' / \texttt{2007.13239v3.pdf}
    \item Leeson and Dwyer, ``Algorithm Selection for Software Verification...'' / \texttt{2201.11711v3.pdf}
  \end{itemize}
\end{frame}

\begin{frame}[standout]
  Read each paper as a surrogate-design case study
\end{frame}

\end{document}
